%% toc: P3 27 27.6,27.7,27.8
\section{Proofs for Keisler's and Shelah's isomorphic-ultrapowers theorems and their hypotheses}
\label{sec:p03-w03-p03-018-annex}

Cardinal arithmetic below uses \cref{thm:p06-cardinal-arithmetic-laws} in the
ambient set theory with Choice.  Two elementary facts are used repeatedly and
are proved here.  For sets \(A\), \(B\), \(C\) write \(B^A\) for the set of
functions from \(A\) to \(B\).
\begin{itemize}
  \item \emph{Monotonicity in the exponent.}  If \(\lambda\ge1\) and
        \(\nu\le\nu'\), then \(\lambda^\nu\le\lambda^{\nu'}\).  Fix
        \(c\in\lambda\) and regard \(\nu\subseteq\nu'\).  Sending
        \(f\in\lambda^\nu\) to its extension by the constant value \(c\) on
        \(\nu'\setminus\nu\) is injective, because the restriction to \(\nu\)
        recovers \(f\).
  \item \emph{Monotonicity in the base.}  If \(\kappa\le\lambda\), then
        \(\kappa^\nu\le\lambda^\nu\), since \(\kappa^\nu\subseteq\lambda^\nu\)
        when \(\kappa\subseteq\lambda\).
\end{itemize}

\begin{annexitem}{annex:p03-w03-p03-018-mu}{thm:p03-w03-p03-018-mu}{The bound \(\mu\) attached to \(\lambda\)}
\begin{proof}
Let \(\lambda\) be infinite.

\emph{Item 1.}  By the base monotonicity fact,
\(2^\lambda\le\lambda^\lambda\), and
\(\lambda^\lambda\le(2^\lambda)^\lambda=2^{\lambda\cdot\lambda}=2^\lambda\)
because \(\lambda<2^\lambda\) and \(\lambda\cdot\lambda=\lambda\).  Hence
\(\lambda^\lambda=2^\lambda>\lambda\).  The class of cardinals \(\nu\) with
\(\lambda^\nu>\lambda\) is therefore nonempty, and as cardinals are well
ordered it has a least element; this is \(\mu\), and \(\mu\le\lambda\).  For
finite \(n\), \(\lambda^0=1\le\lambda\), and \(\lambda^n=\lambda\) for
\(n\ge1\) by induction on \(n\), since
\(\lambda^{n+1}=\lambda^n\cdot\lambda=\lambda\cdot\lambda=\lambda\).  So no
finite cardinal belongs to the class, and \(\mu\) is infinite.

\emph{Item 2.}  Let \(1\le\nu<\mu\).  By minimality of \(\mu\),
\(\lambda^\nu\le\lambda\).  By exponent monotonicity,
\(\lambda=\lambda^1\le\lambda^\nu\).  Hence \(\lambda^\nu=\lambda\).

\emph{Item 3.}  If \(\lambda=\aleph_0\), then
\(\lambda^{\aleph_0}=2^{\aleph_0}>\aleph_0\) by item~1, so
\(\mu\le\aleph_0\); since \(\mu\) is infinite, \(\mu=\aleph_0\).

\emph{Item 4.}  Let \(\lambda=2^\kappa\) with \(\kappa\) infinite.  Then
\(\lambda^\kappa=(2^\kappa)^\kappa=2^{\kappa\cdot\kappa}=2^\kappa=\lambda\).
If \(\mu\le\kappa\), exponent monotonicity would give
\(\lambda<\lambda^\mu\le\lambda^\kappa=\lambda\), which is impossible.  Hence
\(\kappa<\mu\).  Taking \(\kappa=\aleph_0\) gives \(\aleph_0<\mu\) for
\(\lambda=2^{\aleph_0}\), and taking \(\kappa=2^{\aleph_0}\) gives
\(2^{\aleph_0}<\mu\) for \(\lambda=2^{2^{\aleph_0}}\).
\end{proof}
\end{annexitem}

\begin{annexitem}{annex:p03-w03-p03-018-index-bound}{thm:p03-w03-p03-018-index-bound}{An explicit index set for elementarily equivalent models}
\begin{proof}
Let \(M\equiv N\), so \(M\) and \(N\) have one language.  Put
\(\kappa=\lVert M\rVert+\lVert N\rVert+\aleph_0\); this is an infinite
cardinal with \(\lVert M\rVert\le\kappa\) and \(\lVert N\rVert\le\kappa\).
Let \(\lambda=2^\kappa\), an infinite cardinal, and let \(\mu\) be the bound
attached to \(\lambda\).  By item~4 of \cref{thm:p03-w03-p03-018-mu},
\(\kappa<\mu\).  Hence \(\lVert M\rVert<\mu\) and \(\lVert N\rVert<\mu\).
Part~1 of \cref{thm:p03-w03-p03-018-shelah}, applied to this \(\lambda\),
supplies an ultrafilter \(D\) over \(\lambda\) for which every pair of
elementarily equivalent models of power below \(\mu\) has isomorphic
\(\lambda\)-th ultrapowers modulo \(D\); in particular
\(M^\lambda/D\cong N^\lambda/D\).  These ultrapowers are ultrapowers of
\(M\) and of \(N\) in the sense of \cref{sec:types-ultraproducts-05}, with
the common index set \(2^\kappa\) and the common ultrafilter \(D\), which is
the stated form of the forward implication of \cref{thm:p03-keisler-shelah}.
The only cited input is part~1 of \cref{thm:p03-w03-p03-018-shelah}.
\end{proof}
\end{annexitem}

\begin{annexitem}{annex:p03-w03-p03-018-comparison}{thm:p03-w03-p03-018-comparison}{Comparison of the three relations}
\begin{proof}
Elements of \(M^I/D\) are classes \([a]_D\) of functions \(a:I\to M\), and
symbols are interpreted as in \cref{sec:types-ultraproducts-05}.

\emph{Item 1.}  Let \(f:M\to N\) be an isomorphism and define
\(F([a]_D)=[f\circ a]_D\) for \(a:I\to M\).  Because \(f\) is injective,
\(\{i:a(i)=b(i)\}=\{i:f(a(i))=f(b(i))\}\) for all \(a,b:I\to M\); hence
\([a]_D=[b]_D\) iff \([f\circ a]_D=[f\circ b]_D\), so \(F\) is well defined
and injective.  For \(b:I\to N\), \(b=f\circ(f^{-1}\circ b)\), so \(F\) is
surjective.  For an \(n\)-ary relation symbol \(R\),
\[
  \{i:M\entails R(a_1(i),\ldots,a_n(i))\}
  =\{i:N\entails R(f(a_1(i)),\ldots,f(a_n(i)))\}
\]
because \(f\) preserves and reflects \(R\); one set lies in \(D\) iff the
other does, so \(F\) preserves and reflects \(R\) in the ultrapowers.  For an
\(n\)-ary function symbol \(g\), the coordinatewise interpretation and
\(f(g^M(\bar x))=g^N(f\bar x)\) give
\(f\circ g^M(a_1,\ldots,a_n)=g^N(f\circ a_1,\ldots,f\circ a_n)\) pointwise,
so \(F\) commutes with \(g\).  For a constant \(c\), the class of the
constant function \(c^M\) is sent to the class of the constant function
\(f(c^M)=c^N\).  Thus \(F\) is an isomorphism \(M^I/D\cong N^I/D\).

\emph{Item 2.}  By \cref{cor:diagonal-ultrapower}, the diagonal maps
\(M\to M^I/D\) and \(N\to N^J/E\) are elementary embeddings, so
\(M\equiv M^I/D\) and \(N\equiv N^J/E\).  By \cref{lem:iso-invariance}, an
isomorphism \(M^I/D\cong N^J/E\) gives \(M^I/D\equiv N^J/E\).  Since
elementary equivalence is agreement on all sentences, it is transitive and
symmetric, and \(M\equiv N\).

\emph{Item 3.}  Let \(D=\{X\subseteq I:k\in X\}\).  Define
\(e([a]_D)=a(k)\).  For \(a,b:I\to M\), \(\{i:a(i)=b(i)\}\in D\) iff
\(a(k)=b(k)\), so \(e\) is well defined and injective.  Every \(m\in M\) is
\(e\) of the class of the constant function with value \(m\), so \(e\) is
surjective.  For a relation symbol \(R\),
\(R([a_1]_D,\ldots,[a_n]_D)\) holds in \(M^I/D\) iff
\(k\in\{i:M\entails R(a_1(i),\ldots,a_n(i))\}\) iff
\(M\entails R(a_1(k),\ldots,a_n(k))\).  For a function symbol \(g\), the
value of the coordinatewise function \(g^M(a_1,\ldots,a_n)\) at \(k\) is
\(g^M(a_1(k),\ldots,a_n(k))\), and a constant is sent to its value at \(k\),
which is \(c^M\).  Hence \(e\) is an isomorphism \(M^I/D\cong M\).

\emph{Item 4.}  Let \(M\) have exactly \(n\) elements,
\(m_1,\ldots,m_n\), listed without repetition, and let \(M\equiv N\).  The
sentence saying that there are exactly \(n\) elements is true in \(M\), so
\(N\) also has exactly \(n\) elements.  Let \(\mathcal B\) be the finite set
of bijections \(M\to N\).

For a finite set \(L_0\) of nonlogical symbols of the signature, let
\(\delta_{L_0}(x_1,\ldots,x_n)\) be the conjunction of the following
formulas: \(x_j\ne x_l\) for \(j<l\); for each relation symbol \(R\in L_0\) of
arity \(r\) and each \((j_1,\ldots,j_r)\in\{1,\ldots,n\}^r\), the formula
\(R(x_{j_1},\ldots,x_{j_r})\) if \(M\entails R(m_{j_1},\ldots,m_{j_r})\) and
its negation otherwise; for each function symbol \(g\in L_0\) of arity \(r\)
and each such tuple, the formula \(g(x_{j_1},\ldots,x_{j_r})=x_j\) for the
unique \(j\) with \(g^M(m_{j_1},\ldots,m_{j_r})=m_j\); and for each constant
\(c\in L_0\) the formula \(c=x_j\) for the unique \(j\) with \(c^M=m_j\).
This is a finite conjunction, and \(M\entails\delta_{L_0}(m_1,\ldots,m_n)\).
Hence \(M\) satisfies \(\exists x_1\cdots\exists x_n\,\delta_{L_0}\), and so
does \(N\).  If \(b_1,\ldots,b_n\) witness it in \(N\), they are distinct, so
\(m_j\mapsto b_j\) is a bijection; by the conjuncts of \(\delta_{L_0}\) it
preserves and reflects every relation symbol of \(L_0\) and commutes with
every function symbol and constant of \(L_0\).  Let \(\mathcal B_{L_0}\) be
the set of bijections in \(\mathcal B\) with these properties for the symbols
in \(L_0\).  We have shown \(\mathcal B_{L_0}\ne\varnothing\).

If \(L_0\subseteq L_1\), then \(\mathcal B_{L_1}\subseteq\mathcal B_{L_0}\).
Choose a finite \(L_0^\ast\) for which \(\mathcal B_{L_0^\ast}\) has the least
number of elements; such an \(L_0^\ast\) exists because these numbers are
positive integers.  For any finite \(L_1\),
\(\mathcal B_{L_0^\ast\cup L_1}\) is a nonempty subset of
\(\mathcal B_{L_0^\ast}\) with at least as many elements, so it equals
\(\mathcal B_{L_0^\ast}\), and therefore
\(\mathcal B_{L_0^\ast}\subseteq\mathcal B_{L_1}\).  Any element of
\(\mathcal B_{L_0^\ast}\) thus lies in \(\mathcal B_{\{s\}}\) for every symbol
\(s\) of the signature, so it is an isomorphism \(M\cong N\).  When the
signature is finite this step is immediate from the case \(L_0\) equal to the
whole signature.

For the final assertion, let \(M\) be finite.  If \(M\equiv N\), then
\(M\cong N\) by what was just proved; if \(M\cong N\), then
\(M^I/D\cong N^I/D\) by item~1; and if \(M^I/D\cong N^I/D\), then
\(M\equiv N\) by item~2.
\end{proof}
\end{annexitem}

\begin{annexitem}{annex:p03-w03-p03-018-separation}{thm:p03-w03-p03-018-separation}{Isomorphic ultrapowers of non-isomorphic models}
\begin{proof}
\emph{The lexicographic order.}  The relation \(<_{\mathrm{lex}}\) on
\(2^\omega\) is irreflexive by definition.  It is total: distinct \(s,t\)
have a least \(n\) with \(s(n)\ne t(n)\), and then \(s(n)<t(n)\) or
\(t(n)<s(n)\).  For transitivity let \(s<_{\mathrm{lex}}t\), first differing
at \(n\), and \(t<_{\mathrm{lex}}u\), first differing at \(m\).  If \(n=m\),
then \(s(n)=0\), \(t(n)=1\) and \(t(n)<u(n)\), which is impossible for
values in \(\{0,1\}\).  If \(n<m\), then \(s\) and \(u\) agree below \(n\) and
\(s(n)<t(n)=u(n)\).  If \(m<n\), then \(s\) and \(u\) agree below \(m\) and
\(s(m)=t(m)<u(m)\).  In both cases \(s<_{\mathrm{lex}}u\).

\emph{Item 1.}  The order on \(2^\omega\times\mathbb Q\) is a strict linear
order: comparison first by the strict linear order \(<_{\mathrm{lex}}\) and,
on ties, by the strict linear order of \(\mathbb Q\) is irreflexive,
transitive and total by a case check on whether the first coordinates agree.
It has no endpoints, because \((s,q-1)<(s,q)<(s,q+1)\).  It is dense: if
\((s,q)<(t,r)\) with \(s=t\), then \((s,(q+r)/2)\) lies between them; if
\(s<_{\mathrm{lex}}t\), then \((s,q+1)\) lies between them.  The same
witnesses \(q-1\), \(q+1\) and \((q+r)/2\) show that \((\mathbb Q,<)\) is a
dense strict linear order without endpoints.  So \(N\) and
\(M=(\mathbb Q,<)\) are models of \(T_{\rm dlo}\) of
\cref{def:p03-rn-dlo-data}.  That theory is complete by
\cref{thm:p03-rn-dlo-pattern-truth}, so \(M\equiv N\).

\emph{Item 2.}  \(\abs{\mathbb Q}=\aleph_0\), while
\(\abs{2^\omega\times\mathbb Q}=2^{\aleph_0}\cdot\aleph_0=2^{\aleph_0}>\aleph_0\)
by \cref{thm:p06-cardinal-arithmetic-laws}.  An isomorphism is a bijection,
so \(M\not\cong N\).

\emph{Item 3.}  With \(\lambda=2^{2^{\aleph_0}}\), item~4 of
\cref{thm:p03-w03-p03-018-mu} gives \(2^{\aleph_0}<\mu\).  Both models have
power at most \(2^{\aleph_0}\), hence below \(\mu\), and they are
elementarily equivalent by item~1.  Part~1 of
\cref{thm:p03-w03-p03-018-shelah} gives \(M^\lambda/D\cong N^\lambda/D\) for
the ultrafilter \(D\) it supplies.

\emph{Item 4.}  If \(D\) were principal, item~3 of
\cref{thm:p03-w03-p03-018-comparison} would give
\(M\cong M^\lambda/D\) and \(N\cong N^\lambda/D\), and with item~3 above
\(M\cong N\), contradicting item~2.

\emph{Item 5.}  The diagonal map \(N\to N^\lambda/D\) is injective, since
constant functions with distinct values agree on the empty set, which is not
in \(D\).  Hence
\(\lVert M^\lambda/D\rVert=\lVert N^\lambda/D\rVert\ge2^{\aleph_0}>\aleph_0=\lVert M\rVert\),
and
\(M\not\cong M^\lambda/D\).  By \cref{cor:diagonal-ultrapower} the diagonal
map \(d:M\to M^\lambda/D\) is an elementary embedding, so
\(M\equiv M^\lambda/D\) and the image \(d[M]\) is an elementary substructure
of \(M^\lambda/D\).  The image is countable and the ultrapower is not, so the
substructure is proper.
\end{proof}
\end{annexitem}

\subsection*{Keisler's theorems}

The notation is that of \cref{def:p03-w03-p03-018-keisler-notation}.  Three
facts about \(\overline S=\overline S_\alpha(I)\), for an infinite cardinal
\(\alpha\) and a set \(I\) with \(\varkappa(I)=\alpha\), are used below.
\begin{itemize}
  \item[(F1)] \(\overline S\) is closed under supersets within \(I\) and
        under finite intersections.  If \(I-Y\subseteq K\subseteq I\) with
        \(\varkappa(Y)<\alpha\), then \(I-K\subseteq Y\) has fewer than
        \(\alpha\) elements.  And \((I-Y)\cap(I-Y')=I-(Y\cup Y')\), where
        \(\varkappa(Y\cup Y')\le\varkappa(Y)+\varkappa(Y')<\alpha\) because
        \(\alpha\) is infinite (\cref{thm:p06-cardinal-arithmetic-laws}).
  \item[(F2)] A set \(J\subseteq I\) with \(\varkappa(J)=\alpha\) meets every
        member \(I-Y\) of \(\overline S\): otherwise \(J\subseteq Y\) and
        \(\varkappa(J)<\alpha\).  In particular \(\varnothing\notin\overline S\).
  \item[(F3)] For \(J\subseteq I\), \(J\notin\overline S\) if and only if
        \(\varkappa(I-J)=\alpha\), since \(\varkappa(I-J)\le\alpha\).
\end{itemize}

\begin{annexitem}{annex:p03-w03-p03-018-keisler-disjoint}{thm:p03-w03-p03-018-keisler-disjoint}{Keisler's disjoint refinement}
\begin{proof}
Since \(\alpha\) is infinite, \(\alpha\cdot\alpha=\alpha\)
(\cref{thm:p06-cardinal-arithmetic-laws}), so there is a bijection
\(h:\alpha\to\alpha\times\alpha\).  Fix a choice function on the nonempty
subsets of \(\bigcup_{\zeta<\alpha}X_\zeta\).  Define \(g(\xi)\) for
\(\xi<\alpha\) by transfinite recursion: if \(h(\xi)=(\zeta,\eta)\), let
\(g(\xi)\) be the chosen element of
\(X_\zeta-\{g(\xi'):\xi'<\xi\}\).  This set is nonempty, because
\(\varkappa(X_\zeta)=\alpha\) while \(\{g(\xi'):\xi'<\xi\}\) has at most
\(\varkappa(\xi)<\alpha\) elements, as \(\xi<\alpha\) and \(\alpha\) is a
cardinal.  By construction \(g(\xi)\) differs from every earlier value, so
\(g\) is injective, and \(g(\xi)\in X_\zeta\) whenever
\(h(\xi)=(\zeta,\eta)\).

For \(\zeta<\alpha\) put
\(Y_\zeta=\{g(\xi):\xi<\alpha,\ h(\xi)\in\{\zeta\}\times\alpha\}\).  Then
\(Y_\zeta\subseteq X_\zeta\).  For \(\eta\ne\zeta\) the index sets
\(h^{-1}[\{\eta\}\times\alpha]\) and \(h^{-1}[\{\zeta\}\times\alpha]\) are
disjoint, and \(g\) is injective, so \(Y_\eta\cap Y_\zeta=\varnothing\).
Finally \(g\) maps \(h^{-1}[\{\zeta\}\times\alpha]\) injectively onto
\(Y_\zeta\), and that index set is in bijection with
\(\{\zeta\}\times\alpha\) through \(h\), so \(\varkappa(Y_\zeta)=\alpha\).
\end{proof}
\end{annexitem}

\begin{annexitem}{annex:p03-w03-p03-018-keisler-products}{thm:p03-w03-p03-018-keisler-products}{Keisler's ultraproduct theorem}
\begin{proof}
Write \(\overline S\) for \(\overline S_\alpha(I)\) and \(A_i\), \(B_i\) for
the universes of \(\mathfrak A_i\), \(\mathfrak B_i\).

\emph{(i) implies (ii).}  Let \(D\), \(E\) and the isomorphism be as in (i),
and suppose \(\Phi\) is a sentence with
\(\{i:\Phi\in\mathfrak A_i^*\}\in\overline S\subseteq D\) and
\(\{i:(\neg\Phi)\in\mathfrak B_i^*\}\in\overline S\subseteq E\).  By
\L o\'s's theorem (\cref{thm:los}), \(\Phi\) holds in
\(\prod_i\mathfrak A_i/D\) and \(\neg\Phi\) holds in
\(\prod_i\mathfrak B_i/E\).  Isomorphic structures satisfy the same sentences
(\cref{lem:iso-invariance}), a contradiction.

\emph{(ii) implies (i).}  Assume (ii).

\emph{Expanded languages.}  For an ordinal \(\xi\), let \(L_\xi\) be the
language obtained by adding new unary predicate symbols \(P_\zeta\),
\(\zeta<\xi\); the source writes \(P_{\varrho+\zeta}\).  Put
\(A=\prod_{i\in I}A_i\) and \(B=\prod_{i\in I}B_i\).  For a sequence
\(c=\langle c_\zeta\rangle_{\zeta<\xi}\) of elements of \(A\) and
\(i\in I\), let \((\mathfrak A_i,c(i))\) be the expansion of
\(\mathfrak A_i\) to \(L_\xi\) interpreting \(P_\zeta\) as
\(\{c_\zeta(i)\}\); define \((\mathfrak B_i,d(i))\) likewise for a sequence
\(d\) of elements of \(B\).  For a sentence \(\Phi\) of \(L_\xi\) put
\begin{align*}
  Y_c(\Phi)&=\{i\in I:\Phi\text{ holds in }(\mathfrak A_i,c(i))\},\\
  Z_d(\Phi)&=\{i\in I:\Phi\text{ holds in }(\mathfrak B_i,d(i))\}.
\end{align*}
Then \(Y_c(\neg\Phi)=I-Y_c(\Phi)\) and
\(Y_c(\Phi\wedge\Psi)=Y_c(\Phi)\cap Y_c(\Psi)\), and likewise for \(Z_d\).
Say that \((c,d)\) is \emph{separated} when there is no sentence \(\Phi\) of
\(L_\xi\) with \(Y_c(\Phi)\in\overline S\) and
\(Z_d(\neg\Phi)\in\overline S\).  Replacing \(\Phi\) by \(\neg\Phi\) and
using \(Y_c(\neg\neg\Phi)=Y_c(\Phi)\), this is equivalent to: there is no
\(\Phi\) with \(Z_d(\Phi)\in\overline S\) and
\(Y_c(\neg\Phi)\in\overline S\).  So the notion is symmetric in the two
sides.  For \(\xi=0\), separation is hypothesis~(ii).

\emph{Enumerations.}  Each \(A_i\) is nonempty with
\(\varkappa(A_i)\le\varkappa(2^\alpha)=\alpha^+\), so by
\cref{thm:p06-cardinal-arithmetic-laws} and the hypothesis,
\(\varkappa(A)\le(\alpha^+)^\alpha=\varkappa(2^\alpha)^\alpha
=\varkappa(2^{\alpha\cdot\alpha})=\varkappa(2^\alpha)=\alpha^+\).  The
product \(A\) is nonempty by Choice.  By the well-ordering principle there is
therefore a function \(a\) with domain \(\alpha^+\) and range \(A\), and
likewise \(b\) with domain \(\alpha^+\) and range \(B\).  Every ordinal is
uniquely \(\lambda+n\) with \(\lambda\) zero or a limit ordinal and
\(n<\omega\); below, \(\lambda\) ranges over such ordinals.

\emph{The family \(F\).}  Let \(F\) be the set of pairs \((c,d)\) such that,
for some \(\xi<\alpha^+\), \(c\in A^\xi\) and \(d\in B^\xi\), and:
\begin{enumerate}
  \item \(c_{\lambda+2n}=a_{\lambda+n}\) whenever \(\lambda+2n<\xi\);
  \item \(d_{\lambda+2n+1}=b_{\lambda+n}\) whenever \(\lambda+2n+1<\xi\);
  \item \((c,d)\) is separated.
\end{enumerate}
Order \(F\) by extension of both sequences.  The pair of empty sequences is
in \(F\) and lies below every member.  By the maximal principle \(F\) has a
maximal chain \(G\), which contains the empty pair.  Let \(c\) and \(d\) be
the unions of the first and of the second components of members of \(G\).
Their domains are ordinals, linearly ordered by inclusion, so \(c\) and \(d\)
are sequences with a common domain \(\xi\le\alpha^+\).  Conditions~1 and~2
hold for \((c,d)\), since each instance involves one index below \(\xi\),
hence below the domain of some member of \(G\), where the condition holds.
Separation holds as well: a sentence \(\Phi\) of \(L_\xi\) contains finitely
many symbols \(P_\zeta\), all with \(\zeta\) below the domain \(\xi_e\) of
some \((c^e,d^e)\in G\) (any member if there are none), so \(\Phi\) is a
sentence of \(L_{\xi_e}\), and \(Y_c(\Phi)=Y_{c^e}(\Phi)\),
\(Z_d(\neg\Phi)=Z_{d^e}(\neg\Phi)\) because the expansions agree on the
symbols of \(\Phi\).

\emph{The chain has length \(\alpha^+\).}  Suppose \(\xi<\alpha^+\), so
\((c,d)\in F\), and write \(\xi=\lambda+m\) with \(m<\omega\).

Suppose first \(m=2n\).  Put \(c'=c\cup\{(\xi,a_{\lambda+n})\}\), a
sequence of length \(\xi+1\) satisfying condition~1.  Let \(\Gamma\) be the
set of formulas \(\Phi(v_0)\) of \(L_\xi\) with \(v_0\) as their only free
variable such that
\(Y_{c'}(\exists v_0[P_\xi(v_0)\wedge\Phi(v_0)])\in\overline S\).  The
language \(L_\xi\) has \(\varkappa(\varrho)+\varkappa(\xi)\le\alpha\)
non-logical symbols, since \(\varkappa(\xi)\le\alpha\) for \(\xi<\alpha^+\);
formulas are finite strings, so \(\varkappa(\Gamma)\le\alpha\).  The formula
\(v_0=v_0\) is in \(\Gamma\), since the corresponding set is
\(I=I-\varnothing\in\overline S\).

For \(\Phi\in\Gamma\) put \(X(\Phi)=Z_d(\exists v_0\Phi(v_0))\).  Then
\(\varkappa(X(\Phi))=\alpha\).  Indeed, \(\Phi\) does not contain
\(P_\xi\), so if \(c'_\xi(i)\) satisfies \(\Phi\) in
\((\mathfrak A_i,c'(i))\) it satisfies \(\Phi\) in \((\mathfrak A_i,c(i))\);
hence \(Y_c(\exists v_0\Phi)\supseteq
Y_{c'}(\exists v_0[P_\xi(v_0)\wedge\Phi])\), and
\(Y_c(\exists v_0\Phi)\in\overline S\) by (F1).  If \(X(\Phi)\) had fewer than
\(\alpha\) elements, then \(Z_d(\neg\exists v_0\Phi)=I-X(\Phi)\in\overline
S\), and the sentence \(\exists v_0\Phi\) would contradict the separation of
\((c,d)\).

List \(\Gamma\) as \(\langle\Phi_\eta\rangle_{\eta<\alpha}\), with
repetitions if \(\varkappa(\Gamma)<\alpha\).
\Cref{thm:p03-w03-p03-018-keisler-disjoint} applied to
\(X_\eta=X(\Phi_\eta)\) gives sets \(Y_\eta\).  For \(\Phi\in\Gamma\) let
\(X'(\Phi)=Y_\eta\) for the least \(\eta\) with \(\Phi_\eta=\Phi\).  Distinct
members of \(\Gamma\) thus receive disjoint sets
\(X'(\Phi)\subseteq X(\Phi)\) of cardinality \(\alpha\).  Choose \(f\in B\)
as follows: if \(i\in X'(\Phi)\), which happens for at most one \(\Phi\), let
\(f(i)\) be an element of \(B_i\) satisfying \(\Phi\) in
\((\mathfrak B_i,d(i))\), which exists because \(i\in X(\Phi)\); otherwise
let \(f(i)\) be any element of \(B_i\).  Put \(d'=d\cup\{(\xi,f)\}\).
Condition~2 holds for \(d'\), because \(\xi\) has the form \(\lambda+2n\).

The pair \((c',d')\) is separated.  Suppose \(\Psi\) is a sentence of
\(L_{\xi+1}\) with \(Y_{c'}(\Psi)\in\overline S\) and
\(Z_{d'}(\neg\Psi)\in\overline S\).  Renaming bound variables, assume
\(v_0\) does not occur in \(\Psi\), and let \(\Phi(v_0)\) be the result of
replacing each atomic subformula \(P_\xi(u)\) of \(\Psi\) by \(u=v_0\).  It
is a formula of \(L_\xi\) with only \(v_0\) free.  Since \(P_\xi\) denotes
\(\{c'_\xi(i)\}\) in \((\mathfrak A_i,c'(i))\), an induction on \(\Psi\)
shows that \(\Psi\) holds there exactly when \(c'_\xi(i)\) satisfies \(\Phi\)
in \((\mathfrak A_i,c(i))\), that is, exactly when
\(\exists v_0[P_\xi(v_0)\wedge\Phi(v_0)]\) holds there.  So
\(Y_{c'}(\exists v_0[P_\xi(v_0)\wedge\Phi(v_0)])=Y_{c'}(\Psi)\in\overline
S\), and \(\Phi\in\Gamma\).  In the same way, \(\Psi\) holds in
\((\mathfrak B_i,d'(i))\) exactly when \(f(i)\) satisfies \(\Phi\) in
\((\mathfrak B_i,d(i))\).  For \(i\in X'(\Phi)\) it does, so
\(X'(\Phi)\subseteq I-Z_{d'}(\neg\Psi)\).  The right side has fewer than
\(\alpha\) elements, because \(Z_{d'}(\neg\Psi)\in\overline S\), while
\(\varkappa(X'(\Phi))=\alpha\); a contradiction.

So \((c',d')\in F\), and it properly extends \((c,d)\) and therefore every
member of \(G\).  Then \(G\cup\{(c',d')\}\) is a chain properly containing
\(G\), against maximality.  If \(m=2n+1\), the same argument with the two
sides exchanged, using \(d'=d\cup\{(\xi,b_{\lambda+n})\}\) and the symmetric
form of separation, gives the same contradiction.  Hence \(\xi=\alpha^+\).

\emph{The sequences exhaust the products.}  Let \(\eta<\alpha^+\), with
\(\eta=\lambda+n\).  The infinite cardinal \(\alpha^+\) is a limit ordinal
and \(\lambda<\alpha^+\), so \(\lambda+\omega\le\alpha^+\) and
\(\lambda+2n<\alpha^+\).  By condition~1, \(c_{\lambda+2n}=a_\eta\).  So the
range of \(c\) is \(A\), and by condition~2 the range of \(d\) is \(B\).

\emph{The ultrafilters.}  Write \(Y=Y_c\) and \(Z=Z_d\) for sentences of
\(L_{\alpha^+}\); the pair \((c,d)\) is separated.  Let
\(E_0=\{Z(\Phi):Y(\Phi)\in\overline S\}\).  By (F1) and the rules for
\(\wedge\), \(E_0\) is closed under finite intersections.  Each
\(J=Z(\Phi)\in E_0\) has \(\varkappa(J)=\alpha\): otherwise
\(Z(\neg\Phi)=I-J\in\overline S\) together with \(Y(\Phi)\in\overline S\)
contradicts separation.  So, by (F2), every intersection of finitely many
members of \(E_0\cup\overline S\), which has the form \(J\cap K\) with
\(J\in E_0\cup\{I\}\) and \(K\in\overline S\), is nonempty.  By the
Ultrafilter Theorem \cite[480]{Keisler1961ElementaryClasses} there is an
ultrafilter \(E\) on \(I\) with \(E_0\cup\overline S\subseteq E\).

Let \(D_0=\{Y(\Phi):Z(\Phi)\in E\}\), again closed under finite
intersections.  If \(J=Y(\Phi)\in D_0\), then \(Z(\neg\Phi)=I-Z(\Phi)\notin
E\), so \(Z(\neg\Phi)\notin E_0\); hence \(Y(\neg\Phi)=I-J\notin\overline
S\), and \(\varkappa(J)=\alpha\) by (F3).  As before there is an ultrafilter
\(D\) on \(I\) with \(D_0\cup\overline S\subseteq D\).

For every sentence \(\Phi\) of \(L_{\alpha^+}\), \(Y(\Phi)\in D\) if and only
if \(Z(\Phi)\in E\).  If \(Z(\Phi)\in E\), then \(Y(\Phi)\in D_0\subseteq D\).
If \(Z(\Phi)\notin E\), then \(Z(\neg\Phi)\in E\), so
\(Y(\neg\Phi)\in D\) and \(Y(\Phi)\notin D\).

\emph{The isomorphism.}  For \(\eta,\zeta<\alpha^+\),
\(\{i:c_\eta(i)=c_\zeta(i)\}=Y(\exists v_0[P_\eta(v_0)\wedge
P_\zeta(v_0)])\), and similarly for \(d\) with \(Z\).  Hence
\(c_\eta/D=c_\zeta/D\) if and only if \(d_\eta/E=d_\zeta/E\).  Since the
ranges of \(c\) and \(d\) are \(A\) and \(B\), every element of
\(\prod_iA_i/D\) is some \(c_\zeta/D\) and every element of
\(\prod_iB_i/E\) is some \(d_\zeta/E\).  So \(c_\zeta/D\mapsto d_\zeta/E\)
is a well-defined bijection \(h\) from \(\prod_iA_i/D\) onto
\(\prod_iB_i/E\).  Let \(R\) be one of the predicate symbols, with \(k\)
places, and \(\zeta_1,\ldots,\zeta_k<\alpha^+\).  The set of \(i\) at which
\(R\) holds of \((c_{\zeta_1}(i),\ldots,c_{\zeta_k}(i))\) in
\(\mathfrak A_i\) is \(Y(\Theta)\), where \(\Theta\) is
\(\exists v_1\cdots\exists v_k[P_{\zeta_1}(v_1)\wedge\cdots\wedge
P_{\zeta_k}(v_k)\wedge R(v_1,\ldots,v_k)]\); the corresponding set for the
\(d_{\zeta_j}\) in the \(\mathfrak B_i\) is \(Z(\Theta)\).  By the definition
of the ultraproduct relations (\cref{sec:types-ultraproducts-05}), \(R\) holds
of \((c_{\zeta_1}/D,\ldots,c_{\zeta_k}/D)\) exactly when \(Y(\Theta)\in D\),
and of \((d_{\zeta_1}/E,\ldots,d_{\zeta_k}/E)\) exactly when
\(Z(\Theta)\in E\).  These are equivalent, so \(h\) is an isomorphism, and
\(D\), \(E\) contain \(\overline S\).  The source concludes this last step
through \L o\'s's theorem for the expanded structures and their reducts
\cite[486]{Keisler1961ElementaryClasses}; the check above uses only the
definition of the ultraproduct.
\end{proof}
\end{annexitem}

\begin{annexitem}{annex:p03-w03-p03-018-keisler-powers}{thm:p03-w03-p03-018-keisler-powers}{Keisler's ultrapower theorem}
\begin{proof}
\emph{(i) implies (ii).}  If \(\mathfrak A^\alpha/D\cong\mathfrak
B^\alpha/E\), then \(\mathfrak A\equiv\mathfrak B\) by item~2 of
\cref{thm:p03-w03-p03-018-comparison}, that is,
\(\mathfrak A^*=\mathfrak B^*\).

\emph{(ii) implies (i), with the final sentence.}  Take \(I=\alpha\), so
\(\varkappa(I)=\alpha\), and \(\mathfrak A_i=\mathfrak A\),
\(\mathfrak B_i=\mathfrak B\) for all \(i\).  Condition~(ii) of
\cref{thm:p03-w03-p03-018-keisler-products} holds: if
\(\{i:\Phi\in\mathfrak A^*\}\in\overline S_\alpha(\alpha)\), this set is
nonempty by (F2), so \(\Phi\in\mathfrak A^*=\mathfrak B^*\), and then
\(\{i:(\neg\Phi)\in\mathfrak B^*\}=\varnothing\notin\overline
S_\alpha(\alpha)\).  The remaining hypotheses of that theorem are those
assumed here.  It gives ultrafilters \(D\), \(E\) on \(\alpha\) containing
\(\overline S_\alpha(\alpha)\) with
\(\prod_{i<\alpha}\mathfrak A/D\cong\prod_{i<\alpha}\mathfrak B/E\), and these
ultraproducts of constant systems are the ultrapowers
\(\mathfrak A^\alpha/D\) and \(\mathfrak B^\alpha/E\).
\end{proof}
\end{annexitem}

\begin{annexitem}{annex:p03-w03-p03-018-keisler-gch}{thm:p03-w03-p03-018-keisler-gch}{Keisler's existential form}
\begin{proof}
\emph{Forward.}  Let \(\mathfrak A^*=\mathfrak B^*\), with universes \(A\),
\(B\), and let \(\alpha=\varkappa(\varrho)+\varkappa(A)+\varkappa(B)+\omega\),
an infinite cardinal with \(\alpha\ge\varkappa(\varrho)\).  The generalized
continuum hypothesis at \(\alpha\) gives \(\varkappa(2^\alpha)=\alpha^+\), and
\(\varkappa(A),\varkappa(B)\le\alpha<\varkappa(2^\alpha)\).
\Cref{thm:p03-w03-p03-018-keisler-powers} gives ultrafilters \(D\), \(E\) on
\(\alpha\) with \(\mathfrak A^\alpha/D\cong\mathfrak B^\alpha/E\).  Then
\(\mathfrak A^\alpha/D\) belongs to \(\mathrm{Pow}(\mathfrak A)\), and, being
isomorphic to an ultrapower of \(\mathfrak B\), to
\(\mathrm{Pow}(\mathfrak B)\).

\emph{Backward.}  Let \(\mathfrak C\in\mathrm{Pow}(\mathfrak A)\cap
\mathrm{Pow}(\mathfrak B)\): \(\mathfrak C\cong\mathfrak A^I/D\) and
\(\mathfrak C\cong\mathfrak B^J/E\) for ultrafilters \(D\) on \(I\) and \(E\)
on \(J\); the sets \(I\), \(J\) are nonempty because a proper filter contains
its underlying set and not \(\varnothing\).  Then
\(\mathfrak A^I/D\cong\mathfrak B^J/E\), and item~2 of
\cref{thm:p03-w03-p03-018-comparison} gives \(\mathfrak A\equiv\mathfrak B\).
This direction uses no cardinal hypothesis; the source obtains it from the
closure of elementary classes under ultrapowers
\cite[487]{Keisler1961ElementaryClasses}.
\end{proof}
\end{annexitem}

\begin{annexitem}{annex:p03-w03-p03-018-keisler-compare}{thm:p03-w03-p03-018-keisler-compare}{Comparison of the two theorems' hypotheses}
\begin{proof}
\emph{Item 1.}  Suppose \(\overline S_\alpha(I)\subseteq D\) and
\(Y\in D\) with \(\varkappa(Y)<\alpha\).  Then \(I-Y\in\overline
S_\alpha(I)\subseteq D\), so \(\varnothing=Y\cap(I-Y)\in D\), which is
impossible.  Conversely, suppose every member of \(D\) has cardinality
\(\alpha\), and let \(I-Y\in\overline S_\alpha(I)\) with
\(\varkappa(Y)<\alpha\).  Then \(Y\notin D\), so \(I-Y\in D\) because \(D\)
is an ultrafilter.  If \(D\) is principal, generated by \(\{k\}\), then
\(\{k\}\in D\) has cardinality \(1<\alpha\); so a \(D\) of this kind is not
principal.

\emph{Item 2.}  By item~1 of \cref{thm:p03-w03-p03-018-mu} with
\(\lambda=\alpha\), \(\mu\le\alpha\), and \(\alpha<\alpha^+=
\varkappa(2^\alpha)\).  A structure of power less than \(\mu\) therefore has
fewer than \(\varkappa(2^\alpha)\) elements.  The structure with universe
\(\alpha\) and the empty type has \(\alpha\le\varkappa(2^\alpha)\) elements,
and its power \(\alpha\) is at least \(\mu\).

\emph{Item 3.}  If \(\mathrm{Pow}(\mathfrak A)\cap\mathrm{Pow}(\mathfrak
B)\ne\varnothing\), then \(\mathfrak A^*=\mathfrak B^*\) by the backward
direction of the proof of \cref{thm:p03-w03-p03-018-keisler-gch}, which
uses no cardinal hypothesis.  Conversely, if \(\mathfrak A\equiv\mathfrak
B\), \cref{thm:p03-w03-p03-017-keisler-shelah} gives nonempty sets \(I\),
\(J\) and ultrafilters \(U\) on \(I\), \(V\) on \(J\) with
\(\mathfrak A^I/U\cong\mathfrak B^J/V\); this structure lies in
\(\mathrm{Pow}(\mathfrak A)\cap\mathrm{Pow}(\mathfrak B)\).  So the relation
coincides with \(\mathfrak A^*=\mathfrak B^*\), which is an equality of sets
of sentences and hence reflexive, symmetric and transitive.
\end{proof}
\end{annexitem}
